\documentclass[tate,twocolumn,book,paper=a4]{jlreq}

\usepackage{luatexja-fontspec}
\setmainjfont{Noto Serif CJK JP}[
  Path=/usr/share/fonts/opentype/noto/,
  Extension=.ttc,
  UprightFont=NotoSerifCJK-Regular,
  BoldFont=NotoSerifCJK-Bold,
  FontIndex=0,
  BoldFontIndex=0,
]

\title{縦中横・二段見本}
\author{芥川龍之介}
\begin{document}
\maketitle
本書の組版見本は令和8年\tatechuyoko{11}月\tatechuyoko{17}日に作成した。第\tatechuyoko{12}回の試験では\tatechuyoko{ABC}案と\tatechuyoko{12pt}活字を用い、幅は2欄とする。\par
一\par
ある日の事でございます。御釈迦様は極楽の蓮池のふちを、独りでぶらぶら御歩きになっていらっしゃいました。池の中に咲いている蓮の花は、みんな玉のようにまっ白で、そのまん中にある金色の蕊からは、何とも云えない好い匂が、絶間なくあたりへ溢れて居ります。極楽は丁度朝なのでございましょう。\par
やがて御釈迦様はその池のふちに御佇みになって、水の面を蔽っている蓮の葉の間から、ふと下の容子を御覧になりました。この極楽の蓮池の下は、丁度地獄の底に当って居りますから、水晶のような水を透き徹して、三途の河や針の山の景色が、丁度覗き眼鏡を見るように、はっきりと見えるのでございます。\par
するとその地獄の底に、陀多と云う男が一人、ほかの罪人と一しょに蠢いている姿が、御眼に止まりました。この陀多と云う男は、人を殺したり家に火をつけたり、いろいろ悪事を働いた大泥坊でございますが、それでもたった一つ、善い事を致した覚えがございます。と申しますのは、ある時この男が深い林の中を通りますと、小さな蜘蛛が一匹、路ばたを這って行くのが見えました。そこで陀多は早速足を挙げて、踏み殺そうと致しましたが、「いや、いや、これも小さいながら、命のあるものに違いない。その命を無暗にとると云う事は、いくら何でも可哀そうだ。」と、こう急に思い返して、とうとうその蜘蛛を殺さずに助けてやったからでございます。\par
御釈迦様は地獄の容子を御覧になりながら、この陀多には蜘蛛を助けた事があるのを御思い出しになりました。そうしてそれだけの善い事をした報には、出来るなら、この男を地獄から救い出してやろうと御考えになりました。幸い、側を見ますと、翡翠のような色をした蓮の葉の上に、極楽の蜘蛛が一匹、美しい銀色の糸をかけて居ります。御釈迦様はその蜘蛛の糸をそっと御手に御取りになって、玉のような白蓮の間から、遥か下にある地獄の底へ、まっすぐにそれを御下しなさいました。\par
二\par
こちらは地獄の底の血の池で、ほかの罪人と一しょに、浮いたり沈んだりしていた陀多でございます。何しろどちらを見ても、まっ暗で、たまにそのくら暗からぼんやり浮き上っているものがあると思いますと、それは恐しい針の山の針が光るのでございますから、その心細さと云ったらございません。その上あたりは墓の中のようにしんと静まり返って、たまに聞えるものと云っては、ただ罪人がつく微な嘆息ばかりでございます。これはここへ落ちて来るほどの人間は、もうさまざまな地獄の責苦に疲れはてて、泣声を出す力さえなくなっているのでございましょう。ですからさすが大泥坊の陀多も、やはり血の池の血に咽びながら、まるで死にかかった蛙のように、ただもがいてばかり居りました。\par
ところがある時の事でございます。何気なく陀多が頭を挙げて、血の池の空を眺めますと、そのひっそりとした暗の中を、遠い遠い天上から、銀色の蜘蛛の糸が、まるで人目にかかるのを恐れるように、一すじ細く光りながら、するすると自分の上へ垂れて参るのではございませんか。陀多はこれを見ると、思わず手を拍って喜びました。この糸に縋りついて、どこまでものぼって行けば、きっと地獄からぬけ出せるのに相違ございません。いや、うまく行くと、極楽へはいる事さえも出来ましょう。そうすれば、もう針の山へ追い上げられる事もなくなれば、血の池に沈められる事もある筈はございません。\par
こう思いましたから陀多は、早速その蜘蛛の糸を両手でしっかりとつかみながら、一生懸命に上へ上へとたぐりのぼり始めました。元より大泥坊の事でございますから、こう云う事には昔から、慣れ切っているのでございます。\par
（以下、青空文庫本文の抄録）\par
本文は著作権保護期間の満了した作品である。青空文庫の本文を利用した。字体はオープンフォントライセンスの源ノ明朝系書体による。本PDFの組版はクリエイティブ・コモンズ表示四・〇で提供する。\par
\end{document}
